%%%%%%%%%%%%%%%%%%%%%%%%%%%%%%%%%%%%%%%%%%%%%%%%%%%%%%%%%%%%%%%%%%%%%%%%%%%
%
% Generic template for TFC/TFM/TFG/Tesis
%
% $Id: resultados.tex,v 1.7 2016/03/31 10:44:23 macias Exp $
%
% By:
%  + Javier Macías-Guarasa.
%    Departamento de Electrónica
%    Universidad de Alcalá
%  + Roberto Barra-Chicote.
%    Departamento de Ingeniería Electrónica
%    Universidad Politécnica de Madrid
% 
% Based on original sources by Roberto Barra, Manuel Ocaña, Jesús Nuevo, Pedro Revenga, Fernando Herránz and Noelia Hernández. Thanks a lot to all of them, and to the many anonymous contributors found (thanks to google) that provided help in setting all this up.
%
% See also the additionalContributors.txt file to check the name of additional contributors to this work.
%
% If you think you can add pieces of relevant/useful examples, improvements, please contact us at (macias@depeca.uah.es)
%
% You can freely use this template and please contribute with comments or suggestions!!!
%
%%%%%%%%%%%%%%%%%%%%%%%%%%%%%%%%%%%%%%%%%%%%%%%%%%%%%%%%%%%%%%%%%%%%%%%%%%%

\chapter{Resultados}
\label{cha:resultados}

\section{Introducción}
\label{sec:intro-results}

Una vez realizada la implementación de las cuatro metaheurísticas, y decididos todos los hiperparámetros, se puede llevar a cabo
el experimento, en el cual se ejecuta cada una de las instancias un total de cinco veces para cada algoritmo. El objetivo de esta
sección es recoger y exponer los resultados obtenidos de forma estructurada, analizar el rendimiento de cada técnica y contrastar
la calidad de las soluciones con los resultados documentados en el estudio de referencia.

\section{Resultados del experimento}
\label{sec:exp-results}

Mediante un pequeño programa de Python, se pueden correr las diferentes metaheurísticas durante un número determinado de
veces para cada una de las instancias. Una vez se completa una ejecución, el programa recoge la información más relevante para 
el análisis en dos ficheros \textit{csv} auxiliares:
\begin{itemize}
  \item El primer fichero, \textit{results.csv}, almacena los resultados de las distintas ejecuciones.
  \item El segundo fichero, \textit{history.csv}, contiene los resultados de cada una de las iteraciones o generaciones
  individuales de cada algoritmo. Esta información se capta mediante un parámetro denominado \texttt{history} dentro de 
  cada algoritmo. Después de evaluar la mejor solución encontrada en cada iteración del algoritmo, se añade a la variable
  \texttt{history}, y se devuelve junto a la mejor distancia al finalizar la ejecución.
\end{itemize}
Ambos ficheros serán la base de todo el análisis comparativo, ya que, a partir de la información que contienen, se podrá
realizar lo siguiente:
\begin{itemize}
  \item Realizar cálculos estadísticos como la media o desviación típica.
  \item Esbozar gráficas comparativas, tales como curvas de convergencia.
  \item Calcular el \textit{optimalty gap}, o la proximidad de las mejores distancias obtenidas al óptimo global, para obtener
  la calidad de las soluciones proporcionadas por los algoritmos.
\end{itemize}

El fichero \textit{results.csv} presenta el siguiente aspecto:

\begin{verbatim}
  algorithm,instance,execution,time,distance
  GA,berlin52,1,14.2890,8388.0000
  GA,berlin52,2,27.2805,8399.0000
  GA,berlin52,3,10.8290,9497.0000
  GA,berlin52,4,16.2682,8199.0000
  GA,berlin52,5,6.5741,10034.0000
\end{verbatim}

Como se puede apreciar, se almacenan los siguientes campos: metaheurística ejecutada, nombre de la instancia ejecutada, número
de ejecución, tiempo de ejecución en segundos, y mejor distancia encontrada.

Por otro lado, el fichero \textit{history.csv} cuenta con un formato similar:
\begin{verbatim}
  algorithm,instance,execution,iteration,distance
  GA,berlin52,1,1,23575.0000
  GA,berlin52,1,2,23575.0000
  GA,berlin52,1,3,23575.0000
  GA,berlin52,1,4,23575.0000
  GA,berlin52,1,5,23575.0000
  GA,berlin52,1,6,22856.0000
\end{verbatim}

Sin embargo, la información es ligeramente diferente. Aunque también se almacena la metaheurística empleada, la instancia que
se está ejecutando y el número de la ejecución en curso, también se encuentran los valores de la iteración y su mejor distancia
encontrada. En este caso, se puede ver como el algoritmo genético no mejoró su primera mejor solución encontrada para la instancia
\texttt{berlin52} hasta la sexta generación. Este fichero será crucial a la hora de esbozar gráficas de convergencia y 
realizar comparaciones de los valores de \textit{optimality gap} obtenidos.

Mediante un programa diferente de Python, se puede recoger la información de \textit{results.csv} y obtener una tabla de
resultados similar a la del estudio de referencia. Dicha tabla también se exporta a un nuevo fichero auxiliar, 
\textit{summary.csv}, para consultarse y evaluarse posteriormente.

En la tabla~\ref{tab:time} se recoge el tiempo medio de cada instancia por metaheurística. Este valor no se proporciona en el estudio
de referencia, pero sí que cobra importancia en el presente trabajo:

\begin{table}[ht]
    \renewcommand{\arraystretch}{1.3}
    \caption{Tiempo medio de ejecución por metaheurística e instancia.}
    \label{tab:time}
    \begin{center}
        \begin{tabular}{|c|c|c|c|c|}
            \hline
            & \textbf{GA} & \textbf{ACO} & \textbf{SA} & \textbf{TS} \\
            \hline
            \textbf{berlin52} & 15.0482 & 17.2662 & 0.3989 & 3.1461 \\
            \hline
            \textbf{eil51} & 11.8802 & 14.4063 & 0.4602 & 3.9246 \\
            \hline
            \textbf{kroA100} & 28.0709 & 38.4565 & 0.4510 & 19.3140 \\
            \hline
            \textbf{eil101} & 17.9500 & 40.0365 & 0.4052 & 21.1021 \\
            \hline
            \textbf{rat195} & 68.7705 & 161.0066 & 0.5629 & 161.1416 \\
            \hline
            \textbf{d198} & 63.3372 & 161.0186 & 0.9776 & 232.7889 \\
            \hline
        \end{tabular}
    \end{center}
\end{table}

Los resultados muestran una tendencia esperable: para las instancias más pequeñas, con menor número de nodos, el tiempo de
ejecución es menor. Las instancias de tamaño mediano muestran un incremento de tiempo, y por último, las instancias grandes
cuentan con el mayor tiempo de ejecución en todas las metaheurísticas. Sin embargo, un detalle muy llamativo a mencionar es 
la gran velocidad del algoritmo de recocido simulado: incluso con las instancias más grandes, su velocidad de ejecución media no
sobrepasa el segundo. Este dato justifica la ausencia del parámetro \texttt{patience} en este algoritmo, ya que es innecesario: 
el recocido simulado siempre finaliza cuando se alcanza la temperatura final, y además, ocurre extremadamente rápido en 
comparación con las otras técnicas. Por otro lado, el algoritmo de colonia de hormigas tiende a ser el más lento de todos, salvo
en el caso de la instancia \texttt{d198}, para la cual la búsqueda tabú presenta el mayor tiempo de ejecución. Estos valores 
cobran una notable importancia a la hora de comparar los resultados obtenidos de las mejores rutas, ya que un ejor resultado 
puede significar un mayor tiempo de ejecución del algoritmo correspondiente, y por tanto, una mayor y mejor exploración del 
espacio de búsqueda.

En las tablas~\ref{tab:exp_results} y \ref{tab:ref_results} se recogen los resultados obtenidos en el experimento y los 
proporcionados por el estudio de referencia, respectivamente.


\begin{table}[H]
    \renewcommand{\arraystretch}{1.3}
    \caption{Resultados obtenidos tras el experimento.}
    \label{tab:exp_results}
    \begin{center}
        \begin{tabular}{|c|c|c|c|c|c|c|}
            \hline
            \textbf{Instancia} & \textbf{Algoritmo} & \textbf{Media} &
            \textbf{Opt \%} & \textbf{$\sigma$} & \textbf{Mejor} &
            \textbf{Peor} \\
            \hline
            \multirow{4}{*}{berlin52} & GA & 8903.40 & 18.05 & 813.46 & 8199.00 & 10034.00 \\
            \cline{2-7}
            & ACO & 8210.40 & 8.86 & \textbf{197.21} & 8029.00 & 8542.00 \\
            \cline{2-7}
            & SA & 8537.40 & 13.20 & 223.15 & 8288.00 & 8901.00 \\
            \cline{2-7}
            & TS & \textbf{8171.00} & \textbf{8.34} & 241.11 & \textbf{7826.00} & 8452.00 \\
            \hline
            \multirow{4}{*}{eil51} & GA & 470.40 & 10.42 & 22.31 & 453.00 & 509.00 \\
            \cline{2-7}
            & ACO & 481.40 & 13.00 & 8.71 & 470.00 & 488.00 \\
            \cline{2-7}
            & SA & 454.00 & 6.57 & 8.34 & 444.00 & 464.00 \\
            \cline{2-7}
            & TS & \textbf{441.80} & \textbf{3.71} & \textbf{3.42} & \textbf{438.00} & 447.00 \\
            \hline
            \multirow{4}{*}{kroA100} & GA & 41608.40 & 95.51 & 17196.10 & 24551.00 & 61210.00 \\
            \cline{2-7}
            & ACO & 26504.00 & 24.54 & \textbf{513.32} & 25909.00 & 27140.00 \\
            \cline{2-7}
            & SA & 26857.00 & 26.20 & 918.93 & 25562.00 & 27878.00 \\
            \cline{2-7}
            & TS & \textbf{22313.20} & \textbf{4.85} & 600.29 & \textbf{21830.00} & 23331.00 \\
            \hline
            \multirow{4}{*}{eil101} & GA & 1219.20 & 93.83 & 359.44 & 826.00 & 1699.00\\
            \cline{2-7}
            & ACO & 792.00 & 25.91 & 15.23 & 776.00 & 816.00 \\
            \cline{2-7}
            & SA & 878.40 & 39.65 & 59.27 & 815.00 & 951.00 \\
            \cline{2-7}
            & TS & \textbf{668.80} & \textbf{6.33} & \textbf{13.55} & \textbf{645.00} & 678.00 \\
            \hline
            \multirow{4}{*}{rat195} & GA & 6345.20 & 173.15 & 968.18 & 4781.00 & 7162.00 \\
            \cline{2-7}
            & ACO & 3018.40 & 29.94 & 66.83 & 2957.00 & 3114.00 \\
            \cline{2-7}
            & SA & 5043.20 & 117.10 & 219.66 & 4820.00 & 5333.00 \\
            \cline{2-7}
            & TS & \textbf{2573.60} & \textbf{10.79} & \textbf{38.06} & \textbf{2515.00} & 2609.00 \\
            \hline
            \multirow{4}{*}{d198} & GA & 38968.80 & 146.95 & 4815.63 & 32962.00 & 46354.00 \\
            \cline{2-7}
            & ACO & 19552.40 & 23.91 & 598.81 & 18784.00 & 20232.00 \\
            \cline{2-7}
            & SA & 30303.40 & 92.04 & 2080.52 & 27866.00 & 32800.00 \\
            \cline{2-7}
            & TS & \textbf{16348.60} & \textbf{3.60} & \textbf{147.58} & \textbf{16182.00} & 16570.00 \\
            \hline
        \end{tabular}
    \end{center}
\end{table}

\begin{table}[H]
    \renewcommand{\arraystretch}{1.3}
    \caption{Resultados del estudio de referencia.}
    \label{tab:ref_results}
    \begin{center}
        \begin{tabular}{|c|c|c|c|c|c|c|}
            \hline
            \textbf{Instancia} & \textbf{Algoritmo} & \textbf{Media} &
            \textbf{Opt \%} & \textbf{$\sigma$} & \textbf{Mejor} &
            \textbf{Peor} \\
            \hline
            \multirow{4}{*}{berlin52} & GA & 7946.40 & 5.36 & 280.66 & 7546.00 & 8269.00 \\
            \cline{2-7}
            & ACO & \textbf{7922.32} & 5.04 & \textbf{44.55} & 7872.59 & 7985.41 \\
            \cline{2-7}
            & SA & 7960.67 & 5.55 & 44.69 & 7903.77 & 8020.72 \\
            \cline{2-7}
            & TS & 7740.10 & \textbf{2.63} & 148.90 & \textbf{7544.37} & 7937.87 \\
            \hline
            \multirow{4}{*}{eil51} & GA & 454.10 & 6.60 & 1.35 & 452.90 & 455.90 \\
            \cline{2-7}
            & ACO & 467.46 & 9.73 & \textbf{0.91} & 466.54 & 468.43 \\
            \cline{2-7}
            & SA & 439.13 & \textbf{3.08} & 2.29 & 437.42 & 443.04 \\
            \cline{2-7}
            & TS & \textbf{439.10} & \textbf{3.08} & 4.00 & \textbf{434.01} & 443.58 \\
            \hline
            \multirow{4}{*}{kroA100} & GA & 22,726.20 & 6.79 & 504.18 & 22,278.00 & 23,368.00 \\
            \cline{2-7}
            & ACO & 22,941.68 & 7.80 & \textbf{29.83} & 22,908.97 & 22,990.15 \\
            \cline{2-7}
            & SA & \textbf{22,277.50} & \textbf{4.68} & 708.29 & \textbf{21,837.88} & 22,782.68 \\
            \cline{2-7}
            & TS & 22,521.64 & 5.82 & 215.30 & 22,293.45 & 22,794.73 \\
            \hline
            \multirow{4}{*}{eil101} & GA & 685.89 & 9.04 & 3.81 & 680.67 & 689.56\\
            \cline{2-7}
            & ACO & 752.91 & 19.70 & \textbf{3.66} & 748.03 & 757.40 \\
            \cline{2-7}
            & SA & 672.13 & 6.86 & 5.93 & 664.29 & 679.72 \\
            \cline{2-7}
            & TS & \textbf{667.61} & \textbf{6.14} & 5.16 & \textbf{661.66} & 674.41 \\
            \hline
            \multirow{4}{*}{rat195} & GA & 2414.52 & 3.94 & \textbf{5.90} & 2407.45 & 2420.54 \\
            \cline{2-7}
            & ACO & 2465.11 & 6.12 & 40.07 & 2401.43 & 2499.44 \\
            \cline{2-7}
            & SA & 2537.99 & 9.25 & 24.82 & 2497.54 & 2560.45 \\
            \cline{2-7}
            & TS & \textbf{2373.94} & \textbf{2.19} & 11.66 & \textbf{2359.36} & 2388.40 \\
            \hline
            \multirow{4}{*}{d198} & GA & 16,582.86 & 5.09 & 172.00 & 16,405.87 & 16,829.83 \\
            \cline{2-7}
            & ACO & 18,031.92 & 14.27 & 155.01 & 17,806.47 & 18,206.47 \\
            \cline{2-7}
            & SA & 16,380.49 & 3.81 & 247.95 & \textbf{16,035.16} & 16,728.84 \\
            \cline{2-7}
            & TS & \textbf{16,083.48} & \textbf{1.92} & \textbf{32.19} & 16,043.14 & 6,124.63 \\
            \hline
        \end{tabular}
    \end{center}
\end{table}

Contrastando los valores de ambas tablas, se pueden extraer las siguientes conclusiones:

\begin{itemize}
  \item En cuanto a las instancias pequeñas:
  \begin{itemize}
    \item Para \texttt{berlin52}, TS presenta la mejor distancia y el menor \textit{gap}, mientras que ACO cuenta con el mejor
    $\sigma$ y un buen valor medio en ambas tablas.
    \item Para \texttt{eil51}, TS presenta los mejores resultados en prácticamente todas las categorías de ambas tablas, 
    seguido muy de cerca por SA.
  \end{itemize}
  Para las instancias pequeñas, se puede destacar la búsqueda tabú como la técnica que presenta los mejores valores.

  \item En cuanto a las instancias de tamaño mediano:
  \begin{itemize}
    \item Para \texttt{kroA100}, se pueden observar algunas diferencias: TS destaca como la mejor técnica en el experimento,
    mientras que los autores obtienen unos mejores valores en SA. ACO presenta el mejor valor para $\sigma$, notablemente
    superior a las demás técnicas.
    \item Para \texttt{eil101}, TS cuenta con los mejores valores para casi todas las métricas.
  \end{itemize}
  Para las instancias medianas, vuelve a destacar TS como la mejor metaheurística. Por otro lado, GA resalta como la 
  peor técnica en el presente experimento, mientras que es ACO la que cuenta con peores métricas en el experimento de
  referencia.

  \item En cuanto a las instancias grandes:
  \begin{itemize}
    \item Para \texttt{rat195}, TS presenta los mejores resultados en todas las categorías salvo por el valor de $\sigma$ en el 
    estudio de referencia, donde GA es el ganador. Aun así, es GA quien presenta unas métricas muy inferiores respecto al 
    resto de técnicas en este experimento.
    \item Para \texttt{d198}, TS sigue destacando con unas métricas muy superiores. Sin embargo, ACO resalta como la peor técnica
    en el estudio de referencia, pero en el presente experimento es la segunda mejor metaheurística.
  \end{itemize}

  \item Para las instancias grandes, nuevamente destaca TS como la mejor técnica.
\end{itemize}

Como conclusión general, se puede afirmar que los resultados obtenidos muestran una tendencia considerablemente similar a los del 
estudio de referencia, salvo disparidades puntuales. \textit{Grosso modo}, se puede observar como \textbf{la búsqueda tabú
resulta ser la mejor metaheurística} en ambos casos prácticos. No obstante, es imposible negar la notable diferencia 
cuantiativa entre ambos experimentos. La principal causa de ello es, sin duda, el papel del parámetro \texttt{patience}:
en el experimento realizado, el tiempo medio de ejecución para cada caso parte desde apenas unas décimas de segundo, a un 
máximo de casi 4 minutos. Por otro lado, los autores señalan una limitación en su estudio al comentar que las ejecuciones más
largas duraron hasta casi 19 horas. Es lógico pensar que, si se relajara el valor del parámetro \texttt{patience}, o si se 
eliminara directamente, los algoritmos mejorarían aún más los resultados, reduciendo las métricas de mejor distancia encontrada,
así como la peor distancia, la distancia media, su desviación típica, y el \textit{gap} al óptimo global.

Aun así, es importante destacar como, salvando el caso del algoritmo genético, el cual destaca como la peor técnica en el 
presente experimento, los demás resultados presentan una considerable proximidad a los propuestos por los autores, con un 
tiempo de ejecución drásticamente menor, lo cual es el objetivo principal de las metaheurísticas: se renuncia a encontrar el 
óptimo global mediante una ejecución mucho más rápida

\section{Diagramas y gráficos comparativos}
\label{sec:graphs}

En esta sección, se van a mostrar y discutir una serie de diagramas y gráficos, para comparar diferentes aspectos de las
ejecuciones, las instancias y las metaheurísticas.

\subsection{Gráficas de convergencia por instancia}
\label{subsec:conv_inst}

En este apartado se incluyen seis gráficas que muestran la convergencia de los cuatro algoritmos en cada una de las instancias.
El eje X indica el número de iteraciones o generaciones alcanzadas, mientras que el eje Y muestra la mejor distancia
encontrada.

\begin{figure}[H]
    \centering

    \begin{subfigure}[b]{0.48\textwidth}
        \centering
        \includegraphics[width=\textwidth]{convergencia_berlin52.png}
        \caption{Convergencia de \texttt{berlin52}}
        \label{fig:c_berlin52}
    \end{subfigure}
    \hfill
    \begin{subfigure}[b]{0.48\textwidth}
        \centering
        \includegraphics[width=\textwidth]{convergencia_eil51.png}
        \caption{Convergencia de \texttt{eil51}}
        \label{fig:c_eil51}
    \end{subfigure}

    \vspace{0.5cm}

    \begin{subfigure}[b]{0.48\textwidth}
        \centering
        \includegraphics[width=\textwidth]{convergencia_kroA100.png}
        \caption{Convergencia de \texttt{kroA100}}
        \label{fig:c_kroA100}
    \end{subfigure}
    \hfill
    \begin{subfigure}[b]{0.48\textwidth}
        \centering
        \includegraphics[width=\textwidth]{convergencia_eil101.png}
        \caption{Convergencia de \texttt{eil101}}
        \label{fig:c_eil101}
    \end{subfigure}

    \vspace{0.5cm}

    \begin{subfigure}[b]{0.48\textwidth}
        \centering
        \includegraphics[width=\textwidth]{convergencia_rat195.png}
        \caption{Convergencia de \texttt{rat195}}
        \label{fig:c_rat195}
    \end{subfigure}
    \hfill
    \begin{subfigure}[b]{0.48\textwidth}
        \centering
        \includegraphics[width=\textwidth]{convergencia_d198.png}
        \caption{Convergencia de \texttt{d198}}
        \label{fig:c_d198}
    \end{subfigure}

    \caption{Gráficas de convergencia por instancia.}
    \label{fig:conv_graphs}

\end{figure}

Rápidamente se puede observar como la curva de color verde, correspondiente al recocido simulado, presenta una tendencia
claramente diferente al resto. Esto vuelve a demostrar su funcionamiento interno: puesto que no cuenta con el criterio adicional
de parada por estancamiento, su curva de convergencia recorre muchas más iteraciones en el proceso de enfriamiento. Por el
contrario, los otros algoritmos finalizan cuando alcanzan el umbral de estancamiento, esto es, 25 iteraciones sin mejorar.
No obstante, se puede realizar un análisis de las otras tres técnicas: GA destaca por ser la que se toma un mayor número de 
generaciones en converger, mientras que ACO y TS son más rápidas. Nótese como ACO trabaja con generaciones, es decir, 
múltiples soluciones simultáneas, mientras que TS trabaja con una única solución que se mejora progresivamente, y aun así,
presentan una velocidad de convergencia ciertamente similar. Aun así, TS presenta la mayor velocidad de ejecución, precisamente
por su naturaleza.

\subsection{Gráficas de convergencia normalizada por gap\%}
\label{subsec:conv_gap}

En este apartado se muestran otras cuatro gráficas de convergencia, aunque se sustituye el valor de la mejor distancia
encontrada con la proximidad o \textit{gap} al óptimo en cada una de las metaheurísticas. De este modo, se puede comprobar si
un algoritmo escala bien en calidad al incrementar el tamaño de la instancia.

\begin{figure}[H]
    \centering

    \begin{subfigure}[b]{0.48\textwidth}
        \centering
        \includegraphics[width=\textwidth]{convergencia_gap_GA.png}
        \caption{Convergencia de GA}
        \label{fig:c_ga}
    \end{subfigure}
    \hfill
    \begin{subfigure}[b]{0.48\textwidth}
        \centering
        \includegraphics[width=\textwidth]{convergencia_gap_ACO.png}
        \caption{Convergencia de ACO}
        \label{fig:c_aco}
    \end{subfigure}

    \vspace{0.5cm}

    \begin{subfigure}[b]{0.48\textwidth}
        \centering
        \includegraphics[width=\textwidth]{convergencia_gap_SA.png}
        \caption{Convergencia de SA}
        \label{fig:c_sa}
    \end{subfigure}
    \hfill
    \begin{subfigure}[b]{0.48\textwidth}
        \centering
        \includegraphics[width=\textwidth]{convergencia_gap_TS.png}
        \caption{Convergencia de TS}
        \label{fig:c_ts}
    \end{subfigure}

    \caption{Gráficas de convergencia por gap por metaheurística.}
    \label{fig:conv_metah}

\end{figure}

En este caso, se pueden apreciar unos resultados muy diferentes a los del apartado anterior. Se puede observar como el 
algoritmo genético presenta la peor tendencia de convergencia, ya que, a medida que aumenta el tamaño de la instancia, su 
velocidad de convergencia es peor, y tampoco alcanza el valor de \textit{optimality gap} que presentan las instancias más
pequeñas, así como las otras metaheurísticas. Por otra parte, ACO cuenta con la mayor velocidad de convergencia en este 
aspecto para todas las instancias. TS también obtiene una convergencia notablemente buena, y unos valores de \textit{gap}
muy próximos a cero para todas las instancias.

El caso del recocido simulado es particularmente interesante. Si bien las curvas son muy diferentes entre sí, acaban 
convergiendo en valores relativamente próximos a cero. De nuevo, se puede apreciar como el comportamiento de este algoritmo
es, por diseño, claramente diferente a los demás.

\subsection{Diagramas por instancia}
\label{subsec:boxplot_inst}

En este apartado, se incluyen seis diagramas de caja y bigotes, uno por instancia, donde se evalúa la robustez de cada 
algoritmo según su \textit{gap} al óptimo.

\begin{figure}[H]
    \centering

    \begin{subfigure}[b]{0.48\textwidth}
        \centering
        \includegraphics[width=\textwidth]{boxplot_gap_berlin52.png}
        \caption{Diagrama de caja y bigotes para \texttt{berlin52}}
        \label{fig:bp_berlin52}
    \end{subfigure}
    \hfill
    \begin{subfigure}[b]{0.48\textwidth}
        \centering
        \includegraphics[width=\textwidth]{boxplot_gap_eil51.png}
        \caption{Diagrama de caja y bigotes para \texttt{eil51}}
        \label{fig:bp_eil51}
    \end{subfigure}

    \vspace{0.5cm}

    \begin{subfigure}[b]{0.48\textwidth}
        \centering
        \includegraphics[width=\textwidth]{boxplot_gap_kroA100.png}
        \caption{Diagrama de caja y bigotes para \texttt{kroA100}}
        \label{fig:bp_kroA100}
    \end{subfigure}
    \hfill
    \begin{subfigure}[b]{0.48\textwidth}
        \centering
        \includegraphics[width=\textwidth]{boxplot_gap_eil101.png}
        \caption{Diagrama de caja y bigotes para \texttt{eil101}}
        \label{fig:bp_eil101}
    \end{subfigure}

    \vspace{0.5cm}

    \begin{subfigure}[b]{0.48\textwidth}
        \centering
        \includegraphics[width=\textwidth]{boxplot_gap_rat195.png}
        \caption{Diagrama de caja y bigotes para \texttt{rat195}}
        \label{fig:bp_rat195}
    \end{subfigure}
    \hfill
    \begin{subfigure}[b]{0.48\textwidth}
        \centering
        \includegraphics[width=\textwidth]{boxplot_gap_d198.png}
        \caption{Diagrama de caja y bigotes para \texttt{d198}}
        \label{fig:bp_d198}
    \end{subfigure}

    \caption{Diagramas de caja y bigotes por instancia.}
    \label{fig:bp_size}

\end{figure}

De estos diagramas se puede extraer lo siguiente:

Para \texttt{berlin52}, GA obtiene una mediana baja pero caja y bigotes estirados hacia arriba, lo que significa que una
de las ejecuciones salió muy mal. En efecto, si se consulta el fichero auxiliar \textit{results.csv}, se puede observar que 
la ejecución número 5 ha devuelto un valor óptimo mucho peor que las otras. Incluso la ejecución 3 es bastante diferente a 
las iteraciones 1, 2 y 5:

\begin{verbatim}
  GA,berlin52,1,14.2890,8388.0000
  GA,berlin52,2,27.2805,8399.0000
  GA,berlin52,3,10.8290,9497.0000
  GA,berlin52,4,16.2682,8199.0000
  GA,berlin52,5,6.5741,10034.0000
\end{verbatim}

ACO es la técnica más compacta y baja, y TS es ligeramente más ancha pero con un valor de mediana similar. Se observa, por 
tanto, que ACO es más fiable aunque su media final sea ligeramente peor que TS. Por último, SA tiene la caja colapsada y dos
\textit{outliers} simétricos, lo cual es una situación parecida a la de GA: tres ejecuciones coincidieron, y dos se desviaron.
Revisando el fichero auxiliar, se ve como los \textit{outliers} resultan ser la ejecución 2 y 5:

\begin{verbatim}
  SA,berlin52,1,0.4145,8493.0000
  SA,berlin52,2,0.3884,8901.0000
  SA,berlin52,3,0.3998,8489.0000
  SA,berlin52,4,0.3995,8516.0000
  SA,berlin52,5,0.3924,8288.0000
\end{verbatim}

En el caso de \texttt{eil51}, a diferencia del resto de instancias, ACO presenta tanto una mediana como una dispersión superiores
a las de GA, lo que sugiere que su comportamiento no es uniformemente robusto frente a todas las instancias de tamaño pequeño.
GA presenta una caja bastante más baja, aunque cuenta con un outlier muy disperso, lo que significa que solamente una ejecución
se dispersó mucho de la media. En efecto, la culpable es la ejecución 2:

\begin{verbatim}
  GA,eil51,1,12.5354,461.0000
  GA,eil51,2,8.3717,509.0000
  GA,eil51,3,14.2281,453.0000
  GA,eil51,4,11.8213,469.0000
  GA,eil51,5,12.4446,460.0000
\end{verbatim}

SA y TS sí mantienen un aspecto esperable, siendo TS el algoritmo con una caja más compacta y baja.

En \texttt{kroA100} se muestra un patrón extremo: GA tiene una caja descomunal y sin outliers, lo que indica que la ejecución
ha sido errática en esta instancia. ACO y SA muestran una robustez muy similar, prácticamente indistinguible, y TS vuelve a 
ser la más baja y estrecha de nuevo, con un \textit{outlier} moderado que no compromete la lectura general de la técnica.

La instancia \texttt{eil101} muestra un comportamiento similar a la anterior, donde GA presenta una caja muy grande y alargada.
La caja de ACO es muy compacta también, pero la de SA es ligeramente más grande. Y TS es prácticamente lo opuesto a GA, ya que
prácticamente se aprecia una línea plana y muy baja, con un único \textit{outlier} más bajo. Esta dispersión de resultados es
mínima respecto a GA, con lo cual se evidencia una dependencia fuerte de la inicialización aleatoria.

Para \texttt{rat195}, los cuatro algoritmos suben en magnitud, al situarse todas las cajas algo más arriba que en instancias
anteriores. Sin embargo, la robustez se mantiene en lo relativo: GA presenta la caja más ancha, de nuevo sin \textit{outliers},
demostrando una variabilidad en el algoritmo en sí. ACO mantiene el patrón de caja estrecha, y SA cuenta con una caja de
anchura similar a la instancia anterior. TS vuelve a ser, con diferencia, la más consistente, con una caja muy pequeña y baja.

Por último, la instancia \texttt{d198} sí que presenta cambios en GA respecto a las anteriores: la caja es mucho más pequeña,
pero con dos \textit{outliers} muy alejados, así que no se puede realizar ningún consenso sobre cómo de bien o mal van a 
salir las ejecuciones. ACO mantiene su caja compacta y muy por debajo de GA, SA también cuenta con una caja relativamente
estrecha y en un lugar similar a la instancia anterior, y TS vuelve a ser una línea plana muy próxima a cero.

Con todo esto, se refuerza lo anterior: TS no solo gana en media, sino que es el único algoritmo cuya fiabilidad no se degrada
al crecer la instancia. Por el contrario, GA y SA se hacen cada vez menos consistentes según crece el problema.

\subsection{Diagrama de categoría por tamaño}
\label{subsec:boxplot_cat}

En este apartado se incluye un único diagrama que ilustra el \textit{optimality gap} de cada metaheurística, clasificándolo
según los tres diferentes tamaños de instancias:

\begin{figure}[H]
  \centering 
  {\includegraphics[width=0.8\linewidth]{boxplot_categoria_tamano.png}}
  \caption{Diagrama de caja y bigotes por categoría.}
  \label{fig:bp_cat}
\end{figure}

\subsection{Diagrama de tiempo vs tamaño}
\label{subsec:time_size}


%%% Local Variables:
%%% TeX-master: "../book"
%%% End:
